\bigskip\noindent\textbf{Solução 56.}\quad
Trata-se de um problema de minimização com um \emph{vínculo integral}, o que
sugere o método dos multiplicadores de Lagrange. Procuramos um ponto crítico do
funcional
\[
    J[y]=\int_0^L (y')^2\,\d x
\]
na classe das funções $y\in C^1([0,L])$ com $y(0)=y(L)=0$ e
$\int_0^L y\,\d x=A$. Introduzindo um multiplicador $\lambda\in\R$ para o
vínculo, formamos o lagrangiano auxiliar
\[
    F(x,y,y')=(y')^2-\lambda\,y,
\]
cuja equação de Euler--Lagrange $\dfrac{\d}{\d x}F_{y'}-F_y=0$ fornece
\[
    \frac{\d}{\d x}\bigl(2y'\bigr)-(-\lambda)=0
    \quad\Longrightarrow\quad
    2y''+\lambda=0,
    \qquad\text{isto é}\qquad y''=-\frac{\lambda}{2}.
\]
Como $y''$ é constante, $y$ é necessariamente um polinômio quadrático:
\[
    y(x)=-\frac{\lambda}{4}\,x^2+bx+c .
\]
Impomos agora as três condições. De $y(0)=0$ vem $c=0$. De $y(L)=0$ vem
$-\frac{\lambda}{4}L^2+bL=0$, ou seja $b=\frac{\lambda}{4}L$, de modo que
\[
    y(x)=-\frac{\lambda}{4}\,x^2+\frac{\lambda}{4}L\,x
        =\frac{\lambda}{4}\,x(L-x).
\]
Resta usar o vínculo. Como $\int_0^L x(L-x)\,\d x=\frac{L^3}{6}$, obtemos
\[
    A=\int_0^L y\,\d x=\frac{\lambda}{4}\cdot\frac{L^3}{6}=\frac{\lambda L^3}{24},
    \qquad\text{donde}\qquad \lambda=\frac{24A}{L^3}.
\]
Substituindo, $\frac{\lambda}{4}=\frac{6A}{L^3}$ e portanto
\[
    \boxed{\,y(x)=\frac{6A}{L^3}\,x(L-x)\,}.
\]
Como verificação do vínculo,
$\int_0^L \frac{6A}{L^3}x(L-x)\,\d x=\frac{6A}{L^3}\cdot\frac{L^3}{6}=A$.
Que se trata de fato de um \emph{mínimo} é claro: o integrando $(y')^2$ é convexo
em $y'$, e o conjunto das funções admissíveis (definido por restrições afins) é
convexo; logo $J$ é convexo nessa classe e o único ponto crítico é o minimizador
global.

\bigskip\noindent\textbf{Solução 57.}\quad
O integrando do tempo de descida,
\[
    F(y,y')=\frac{\sqrt{1+y'^2}}{\sqrt{y}},
\]
\emph{não depende explicitamente da variável independente $x$}. Essa é
exatamente a situação em que vale a \emph{identidade de Beltrami}: ao longo de
todo extremal a quantidade $F-y'\,F_{y'}$ é constante.

\begin{enumerate}[label=(\alph*)]
\item Recordemos a identidade de Beltrami. Para um lagrangiano $F=F(y,y')$, a
equação de Euler--Lagrange é $\frac{\d}{\d x}F_{y'}=F_y$, e então
\[
    \frac{\d}{\d x}\bigl(F-y'F_{y'}\bigr)
    =F_y\,y'+F_{y'}\,y''-y''F_{y'}-y'\frac{\d}{\d x}F_{y'}
    =y'\Bigl(F_y-\frac{\d}{\d x}F_{y'}\Bigr)=0 .
\]
Logo $F-y'F_{y'}=\text{const}$. Calculemos. Temos
\[
    F_{y'}=\frac{1}{\sqrt{y}}\cdot\frac{y'}{\sqrt{1+y'^2}},
\]
de modo que
\[
    F-y'F_{y'}
    =\frac{\sqrt{1+y'^2}}{\sqrt{y}}-\frac{y'^2}{\sqrt{y}\,\sqrt{1+y'^2}}
    =\frac{(1+y'^2)-y'^2}{\sqrt{y}\,\sqrt{1+y'^2}}
    =\frac{1}{\sqrt{y}\,\sqrt{1+y'^2}}.
\]
Igualando a uma constante e elevando ao quadrado obtemos
$y\,(1+y'^2)=\text{const}$. Como $y>0$ e a constante é o inverso do quadrado da
constante de Beltrami, ela é positiva; chamando-a de $a>0$,
\[
    \boxed{\,y\bigl(1+y'^2\bigr)=a\,}.
\]

\item Integremos a relação de primeira ordem $y(1+y'^2)=a$. Dela,
\[
    y'^2=\frac{a-y}{y}
    \quad\Longrightarrow\quad
    \frac{\d y}{\d x}=\sqrt{\frac{a-y}{y}}
    \quad\Longrightarrow\quad
    \d x=\sqrt{\frac{y}{a-y}}\;\d y .
\]
A substituição que racionaliza essa raiz é a parametrização angular
\[
    y=a\sin^2\!\Bigl(\frac{\varphi}{2}\Bigr)=\frac{a}{2}\,(1-\cos\varphi),
\]
para a qual $a-y=a\cos^2(\varphi/2)$ e $\d y=a\sin(\varphi/2)\cos(\varphi/2)\,\d\varphi$.
Logo
\[
    \d x=\sqrt{\frac{\sin^2(\varphi/2)}{\cos^2(\varphi/2)}}\;
         a\sin(\varphi/2)\cos(\varphi/2)\,\d\varphi
        =a\sin^2\!\Bigl(\frac{\varphi}{2}\Bigr)\,\d\varphi
        =\frac{a}{2}\,(1-\cos\varphi)\,\d\varphi,
\]
e integrando (escolhendo a origem de $\varphi$ no ponto inicial),
\[
    x=\frac{a}{2}\,(\varphi-\sin\varphi),
    \qquad
    y=\frac{a}{2}\,(1-\cos\varphi).
\]
Pondo $R:=a/2$, isto é precisamente a ciclóide
\[
    \boxed{\;x(\varphi)=R\,(\varphi-\sin\varphi),
        \qquad y(\varphi)=R\,(1-\cos\varphi)\;}
\]
descrita por um ponto da circunferência de raio $R$ que rola sobre o eixo $x$ (a
parametrização do enunciado corresponde a renomear $R$ por $a$). Verifiquemos
diretamente que essa curva satisfaz a equação extremal. De
$\frac{\d y}{\d\varphi}=R\sin\varphi$ e
$\frac{\d x}{\d\varphi}=R(1-\cos\varphi)$ vem
\[
    y'=\frac{\sin\varphi}{1-\cos\varphi},
    \qquad
    1+y'^2=\frac{(1-\cos\varphi)+(1+\cos\varphi)}{1-\cos\varphi}
          =\frac{2}{1-\cos\varphi},
\]
onde usamos $\sin^2\varphi=(1-\cos\varphi)(1+\cos\varphi)$. Portanto
\[
    y\,(1+y'^2)=R(1-\cos\varphi)\cdot\frac{2}{1-\cos\varphi}=2R,
\]
constante, e a identificação $y(1+y'^2)=a$ dá $a=2R$, em acordo com $R=a/2$.

\item A ciclóide construída em (b) tem um parâmetro livre, o raio $R$ (ou
$a=2R$), e parte da origem em $\varphi=0$ com a orientação fixada. Exigir que ela
\emph{passe pelo ponto final} $(x_1,y_1)$ impõe duas equações,
\[
    x_1=R\,(\varphi_1-\sin\varphi_1),
    \qquad
    y_1=R\,(1-\cos\varphi_1),
\]
nas duas incógnitas $R$ e $\varphi_1$ (o valor de $\varphi$ no ponto de chegada).
Tomando o quociente para eliminar $R$,
\[
    \frac{x_1}{y_1}=\frac{\varphi_1-\sin\varphi_1}{1-\cos\varphi_1},
\]
que é uma única equação na incógnita $\varphi_1$. O lado direito é uma função
contínua e estritamente crescente de $\varphi_1\in(0,2\pi)$, indo de $0$ (quando
$\varphi_1\to0^+$) a $+\infty$ (quando $\varphi_1\to2\pi^-$); logo, para qualquer
razão $x_1/y_1>0$, existe um único $\varphi_1$ que a resolve. Determinado
$\varphi_1$, qualquer das duas equações acima fixa $R=\dfrac{y_1}{1-\cos\varphi_1}$.
Assim, o ponto final $(x_1,y_1)$ determina de modo único o raio $R$ (equivalentemente,
a constante $a$) e o intervalo $\varphi\in[0,\varphi_1]$ percorrido pelo arco.
\end{enumerate}

\bigskip\noindent\textbf{Solução 58.}\quad
O funcional $E[y]=\int_0^1 (y'')^2\,\d x$ depende da \emph{segunda} derivada, de
modo que a condição de extremo é a equação de Euler--Poisson, a generalização da
equação de Euler--Lagrange para lagrangianos $F(x,y,y',y'')$:
\[
    F_y-\frac{\d}{\d x}F_{y'}+\frac{\d^2}{\d x^2}F_{y''}=0 .
\]
Aqui $F=(y'')^2$, logo $F_y=F_{y'}=0$ e $F_{y''}=2y''$, e a equação reduz-se a
\[
    \frac{\d^2}{\d x^2}\bigl(2y''\bigr)=0
    \quad\Longrightarrow\quad
    y^{(4)}=0 .
\]
A solução geral de $y^{(4)}=0$ é um polinômio de grau $\le 3$:
\[
    y(x)=ax^3+bx^2+cx+d .
\]
(As quatro condições de contorno do enunciado fixam todos os quatro coeficientes
e são exatamente as condições naturais para um problema de segunda ordem com
$y$ e $y'$ prescritos nos extremos.) Como $y'(x)=3ax^2+2bx+c$, as condições dão
\[
    \begin{aligned}
    y(0)=0:&\quad d=0,\\
    y'(0)=\alpha:&\quad c=\alpha,\\
    y(1)=0:&\quad a+b+c+d=0,\\
    y'(1)=\beta:&\quad 3a+2b+c=\beta.
    \end{aligned}
\]
Das duas primeiras, $d=0$ e $c=\alpha$. As duas últimas tornam-se
\[
    a+b=-\alpha,
    \qquad
    3a+2b=\beta-\alpha .
\]
Subtraindo o dobro da primeira da segunda, $a=(\beta-\alpha)-2(-\alpha)=\alpha+\beta$,
e então $b=-\alpha-a=-2\alpha-\beta$. Portanto
\[
    \boxed{\,y(x)=(\alpha+\beta)\,x^3-(2\alpha+\beta)\,x^2+\alpha\,x\,},
\]
isto é $a=\alpha+\beta$, $b=-(2\alpha+\beta)$, $c=\alpha$, $d=0$. Conferindo as
condições nos extremos: $y(1)=(\alpha+\beta)-(2\alpha+\beta)+\alpha=0$ e
$y'(1)=3(\alpha+\beta)-2(2\alpha+\beta)+\alpha=\beta$, como exigido. Esse cúbico é
de fato o minimizador: o integrando $(y'')^2$ é convexo em $y''$ e o conjunto
admissível é afim, de sorte que $E$ é convexo e o único extremal é o mínimo
global.
